\section{Conclusion}

We present~\method, a self-evolving framework for interpretable alpha mining that formulates factor discovery as a constrained multi-agent research process. Extensive experiments across both Chinese and U.S. equity markets show that \method consistently produces more stable and generalizable factors than all baselines. Beyond empirical performance, \method emphasizes three properties critical for real-world quantitative research: \emph{diversity}, \emph{controllability}, and \emph{trustworthiness}. Diversity is promoted by broad hypothesis exploration and redundancy-aware evolution, improving robustness under non-stationarity; controllability is enabled by symbolic representations and synthesis-time gates that enforce consistency and limit complexity. Trustworthiness is further strengthened by trajectory-level inheritance, where crossover recombines high-performing segments from previously successful trajectories, providing a verifiable lineage of validated mechanisms and repair patterns. Future work will explore extending \method to multi-asset and cross-market settings, incorporating adaptive regime-aware evolution, and integrating the agentic discovery process more tightly with portfolio construction and risk management. More broadly, we believe that agentic evolution provides a promising paradigm for discovery problems in high-noise, non-stationary domains beyond finance.


\section*{Acknowledgments}
This work was supported by the National Social Science Fund of China Project under Grant No. 22BTJ031; and the Shanghai Engineering Research Center of Finance Intelligence under Grant No. 19DZ2254600. We acknowledge the technical support from the Qinghai Provincial Key Laboratory of Big Data in Finance and Artificial Intelligence Application Technology. 